\section{The excursion reduction in detail}\label{app:excursion}
This appendix contains the local algebra used in Theorem~\ref{thm:height}. Its hypotheses are essential: the matrices are actual prefixes of an MA input, and the scalar ceiling is inherited from that input. An arbitrary proper ordering or an arbitrary path between two small endpoints would not suffice.

The disconnected case was disposed of in the proof of Theorem~\ref{thm:height}; throughout this appendix the support is connected.

Throughout $T\ge H(B)$, $F=T+1$, and a state is called \emph{low} if its height is at most $T$, and \emph{high} if its height is $F$. These terms are used only in this appendix. All trajectory states under consideration have height at most $F$. At every actual prefix and in every signed permutation frame, every totally proper scalar $m$ satisfies
\begin{equation}\label{eq:prefixscalar}
 0\le 2|P|\le m\le S(B)\le T.
\end{equation}
For connected support, $m\ge3$ and $2|P|<m$. To obtain the assertion for a supplied totally proper ordering at a prefix, transport it along the reverse of that prefix to the input. Relabelling transports the pivot letters and global arrow reversal reverses the cyclic order. The resulting scalar is one of the six in $S(B)$.

\subsection{A strong first pair cannot return}
\begin{lemma}\label{lem:strong}
Let $P$ be an actual framed prefix. Suppose $\htq(P)=F$, $\htq(\mu_iP)\le T$, $\htq(\mu_jP)\le F$, $i\ne j$, and $|p_{ij}|\ge2$. No reduced word starting with $j$, and staying under the ceiling $F$, can reach a low state. Here reduced means that consecutive letters differ.
\end{lemma}
\begin{proof}
Put $i,j$ in positions $0,1$. Every entry incident to this pair is below $F$. For example, after orienting $p_{01}>0$, an entry $p_{12}=F$ which decreases under mutation at $0$ requires $p_{02}>0$; mutation at $1$ then creates $p_{02}+p_{01}F>F$. The negative wall and the other incident entries follow by reversal or exchange of the two remaining vertices. The sole possible top entry is therefore $23$. Reversal and exchanging $2,3$ put the state in the form
\[
 P=(a,b,-c,d,E,F),\qquad a\ge2,\quad b,c>0.
\]
The low child and the other ceiling give $b<F$, $ac-E<F$, and $b+ad\le F$ when $d\ge0$. A totally proper scalar $m$ exists and satisfies \eqref{eq:prefixscalar}, with Pfaffian $aF-bE-cd$.

We claim $d\ge0$ and $E\ge c$. If $d<0$, then $E>0$ creates $23$ entry $F+(-d)E>F$ under mutation at $1$, while $E\le0$ gives $P\ge2F$. If $d=0,E<0$ the same Pfaffian contradiction applies. The remaining two excluded chambers have a unique proper class (including its wiggles when a coordinate is zero), and its scalar is the following polynomial:
\begin{align*}
 K={}&a^2+b^2+c^2+d^2+E^2+abd-acE
       +F\{F-bc-d(ac-E)\} &&(0\le E<c),\\
 K_+={}&a^2+b^2+c^2+d^2+e^2+abd+ace
       +F\{F-bc-d(e+ac)\} &&(E=-e<0,d>0).
\end{align*}
Here is the arithmetic exclusion, retaining the guards needed in each case. In the first row suppose $0\le K\le F$ and $2|aF-bE-cd|\le K$. Set $x=b+(a-1)d$ and $y=c-E$. The height inequalities give $1\le x,y\le F-1$. The maximum of $x^2+y^2-Fxy$ on this integer rectangle is attained on its boundary, and is at most $2-F$ (fix one variable and use convexity, then repeat). Writing $p=aF-bE-cd$, the exact identity is
\[
 K=a^2+x^2+y^2-Fxy-(a-1)F^2+Fp
 -(a-2)\{d(b+ad)+cE+FdE\}.
\]
The braced term is nonnegative. Applying the rectangle bound and $2p\le F$ gives
\[
 2K\le2a^2+4-2F-(2a-3)F^2<0.
\]
For $a=2$ use $F\ge3$; for $a\ge3$ use $F>a$, which follows since $a$ is an incident entry. In the second row $e+ac<F$ and $b+(a-1)d<F$. Put $p=aF+be-cd$. Direct expansion gives
\[
 K_+=a^2+(b^2+d^2+abd)+(c^2+e^2+ace)
 -(a^2-1)F^2+apF-F(bc+de+abe).
\]
The two quadratic groups are bounded by $(b+(a-1)d)^2$ and $(ac+e)^2$, respectively, hence by $(F-1)^2$ each. Dropping the last nonpositive term and using $2p\le F$ gives
\[
 2K_+\le-(2a^2-a-6)F^2+2a^2-8F+4<0.
\]
For $a=2$ the last expression is $12-8F<0$; for $a\ge3$ its quadratic term is negative. Both contradict \eqref{eq:prefixscalar}. Thus $d\ge0,E\ge c$.

The next state is
\begin{equation}\label{eq:strongnext}
 \mu_1P=(-a,b+ad,-c+aE,-d,-E,F).
\end{equation}
Mutation at $2$ creates $03$ entry $-c+aE+(b+ad)F>F$. Mutation at $3$ exceeds the ceiling unless $d=0,E=c=1$. On that exceptional wall the proper scalar is
$m=a^2+b^2+F^2-a-bF+2>F$; this follows from
$b^2-bF\ge-F^2/4$, $a\ge2$, and $F\ge3$. Therefore reducedness forces the next pivot to be $0$.

The two columns at the remaining vertices are now aligned: in the orientation where the connecting edge is positive, one pivot row is nonpositive, the other nonnegative, and the latter dominates in absolute value. If $C_{i\ell}\le0\le C_{j\ell}$, $-C_{i\ell}\le C_{j\ell}$ and $C_{ij}\ge2$, mutation at $j$ replaces the entries by $C_{i\ell}+C_{ij}C_{j\ell}$ and $-C_{j\ell}$, preserving dominance. After one further pair mutation both signs are strict. The opposite $F$ edge is unchanged. A pivot at either remaining vertex then changes a strictly signed incident entry by at least $F$, violating the ceiling. Induction forces alternation of $0,1$ forever within any allowed continuation. Its opposite edge remains $F$, so it never returns low.
\end{proof}

\subsection{Zero edges and a bounded commuting square}
\begin{lemma}\label{lem:zero}
Under the prefix and ceiling hypotheses of Lemma~\ref{lem:strong}, but with $p_{ij}=0$, a reduced first return from $P$ starting with $j$ is either $[j]$ or $[j,i]$. If a low state $R$ has $r_{ij}=0$ and a word using just $i,j$ has a low endpoint, that same labelled endpoint has a route of length at most two with every state low.
\end{lemma}
\begin{proof}
Normalize $P=(0,b,-c,d,E,F)$ with $b,c>0$. All five incident entries are low, because mutation at one of the zero pair leaves the other's row unchanged. Except for $d<0<E$, which already violates the ceiling at $1$, its proper classes have scalar
\[
 m=b^2+c^2+d^2+E^2+F^2-bcF+dEF.
\]
If mutation at $2$ stays under $F$, its $03$ entry $-c+bF$ forces $b=1$; its other entries give $d\le1$, and $E\le0$ if $d=1$. When $d\le0$ the excluded sign chamber ensures $dE\ge0$, and
\[
 4(m-F)=(2c-F)^2+3(F-1)^2+2F+1+4d^2+4E^2+4dEF>0.
\]
When $d=1$,
\[
 4(m-F)=(2c-F)^2+(2E+F)^2+2(F-1)^2+6>0.
\]
Both contradict $m\le T$. Mutation at $3$ is the reversed and exchanged argument.

If $\mu_1P$ is still high its opposite entry is $F$ or $-F$. The latter requires $d>0,E=-e<0$ and $de=2F$ with $1\le b,c,d,e\le F-1$. The elementary rectangle and factor bounds
\[
 b^2+c^2-Fbc\le2-F,\qquad
 3(d^2+e^2)\le2F^2-5F+36
\]
give $3m\le42-8F-F^2<0$, since the factor wall forces $F\ge4$. Hence the second high state still has opposite entry $F$. Its other-pivot child has opposite entry $F-bc$, the same as the low $\mu_0P$. The preceding calculation seals this second state against a third pivot. Reducedness and first return give the two listed words.

It remains to bound the alternate square. The two mutations commute and are involutions, so there are at most four pair-orbit states. Only the double-mutation endpoint needs attention. If one intermediate state is high, normalize it as
\[
 P=(0,b,-c,d,-e,f),\qquad b,c,d,e>0,\quad f>T.
\]
The two low children have opposite entries $f-bc$ and $f-de$, and all incident entries are at most $T$. The alternative corner is low precisely when
\begin{equation}\label{eq:zerosquare}
 -T\le f-bc-de\le T.
\end{equation}
The upper inequality is immediate. Its scalar is
$m=b^2+c^2+d^2+e^2+f^2-f(bc+de)$, with $3\le m\le T$ and $2|be-cd|<m$.

If all four incident magnitudes are at least two and the lower inequality fails, put $p=bc$, $q=de$, $U=b^2+c^2$, $V=d^2+e^2$, exchanging pairs so that $q\le p$. The low-child and overflow inequalities imply
\[
 4U-p^2\le16,\quad4V\le2pq,\quad q\ge4,\quad p<2f<p+2q.
\]
The identity
\[
 4m=(4U-p^2)+(4V-2pq)-(2f-p)(p+2q-2f)
\]
then yields $4m\le16-(2q-1)\le9$, impossible.

Otherwise the symmetries $(b,c,d,e)\mapsto(c,b,e,d)$ and $(b,c,d,e)\mapsto(d,e,b,c)$ put a unit entry at $b=1$. Put $u=f-c>0$, $v=de-f>0$. Then
\[
 de=c+u+v,\quad c,d,e,u,v\le T,\quad T<c+u,c+v,\quad2u,2v<de,
\]
and
\begin{equation}\label{eq:zeroM}
 m=d^2+e^2+1-de(u+v)+u^2+uv+v^2,\qquad P=e-cd.
\end{equation}
If $c=1$, then $u=v$, $de=1+2u$, $d,e\ge3$ and $m=2+d^2+e^2-(u+1)^2\le1$. For $c\ge2$ the exhaustive comparisons are
\begin{center}\begin{tabular}{ll}\toprule
Condition & Consequence\\\midrule
All $c\ge2$ & $m\le(d-e)^2+4$\\
$e\le d$, $c<d$ & $m\le2$\\
$e\le d\le c$ & $m\le2|e-cd|$\\
$d<e$, $c<e$ & $m\le2$\\
$d<e$, $c=e$ & $m\le4\le2|e-cd|$\\\bottomrule
\end{tabular}\end{center}
For the first comparison, when $u\le v$, the exact difference is
\[
 (de-2v-1)(u+v-2)+(u-1)^2+3(u-1)(v-u)+(v-u)^2\ge0;
\]
exchange $u,v$ otherwise. For the second use
$m=(d-c)^2+e^2+1-cd(e-2)-uv$ and $u,v\ge d-c+1$; the fourth is the exchanged numerical comparison. The third follows from the first and $c\ge d\ge e\ge2$. At $c=e$ use $m=d^2-(d-2)e^2+1-uv\le4$. Thus the only unexcluded range is $2\le d<e<c$. But $f<de<2(cd-e)<m$, contradicting $m\le T<f$. This proves \eqref{eq:zerosquare}. The low square reaches the identical labelled endpoint, preserving the literal remaining tail and removing all high states of this pair block.
\end{proof}

\subsection{Unit pairs and transport of the remaining word}\label{app:unit}
For $|C_{01}|=1$, the unit pentagon identity is
\begin{equation}\label{eq:pentagon}
 \mut{0,1,0,1,0}C=C^{(1,0,2,3)}.
\end{equation}
It holds for arbitrary signs of the other entries. One can verify each outside column by the alternating positive-part transformations in \eqref{eq:mutation}; the $23$ entry is then forced by $P(\mu_kC)=-P(C)$ and the final unit $01$ entry. Thus only finitely many sign choices in two column variables are needed, with no multiplicity cutoff. Reversal covers the negative-unit case.

For clarity about replacement inside a longer word, suppose $\mu_vC=(\mu_uC)^p$. Then for every chronological tail $w$,
\begin{equation}\label{eq:tailtransport}
 \mu_{v\,p^{-1}(w)}C=(\mu_{uw}C)^p.
\end{equation}
This is proved one letter at a time from permutation covariance of mutation. Corresponding tail states have identical heights and the same high-state count; acyclicity is unchanged. In the original labels, conjugate $p$ by the frame used for normalization.

\begin{lemma}[Bounded unit-pair routes]\label{lem:unitroute}
Let $C$ be MA and $|C_{ij}|=1$. If $C$ and the endpoint of a word using only $i,j$ have height at most $K$, then a word of length at most three, staying in that box, reaches the same endpoint up to the swap of $i,j$.
\end{lemma}
\begin{proof}
Normalize the pair to $01$ and set $N_r=\mu_{u_r}C$ for
$(u_0,u_1,u_2,u_3,u_4)=(\varnothing,0,01,010,0101)$ and $\sigma=(1,0,2,3)$. The ten labelled states are $N_r,N_r^\sigma$. Their first five transition rows are
\begin{center}\begin{tabular}{c|cc}
 & $0$ & $1$\\\hline
$N_0$ & $N_1$ & $N_4^\sigma$\\
$N_1$ & $N_0$ & $N_2$\\
$N_2$ & $N_3$ & $N_1$\\
$N_3$ & $N_2$ & $N_4$\\
$N_4$ & $N_0^\sigma$ & $N_3$
\end{tabular}\end{center}
The other rows follow by $\sigma$, exchanging the pivot labels. If the endpoint is $N_2$, use $01$ when $N_1$ is low and $101$ otherwise; for $N_3$ use $10$ through $N_4$ when low and $010$ otherwise, all up to $\sigma$. Endpoints $N_0,N_1,N_4$ need at most one mutation.

The nontrivial bound is the alternate arc past a strict unit peak. Normalize such a peak as $P=(1,b,-c,d,-e,f)$, where $b,c,d,e\ge1$, $f>K$, and its two children are in the height-$K$ box. Then $b+d\le K$, $c+e\le K$. Set
\[
 D=bc+de+cd,\quad m=1+b^2+c^2+d^2+e^2+bd+ce+f^2-fD,\quad p=f+be-cd.
\]
The alternate arc's two interior tuples are
\[
 (1,-b,-e,-d,c+e,f-D),\quad(-1,b,e,-b-d,c,f-D).
\]
Thus its only remaining guard is $f-D\ge-K$. If it fails, $f+K<D$. Expansion, with the displayed low-child bounds, gives
\begin{center}\begin{tabular}{ll}\toprule
Range & Contradiction to $0\le2|p|\le m$\\\midrule
$c=1$ or $d=1$ & $p>0$, $m-2p\le-5$\\
$c,d\ge2$, $b=1$ or $e=1$ & $m\le-1$\\
$b,c,d,e\ge2$ & $m\le-10$\\\bottomrule
\end{tabular}\end{center}
The products giving these bounds can be written explicitly. In the last row use
$m=1+(b+d)^2+(c+e)^2-bd-ce-f(D-f)$,
$D\ge2(b+d)+2(c+e)-4$ and
$(f-X-1)(D-f-X-1)\ge0$, where $X=\max(b+d,c+e)$.
If $b=1$, put $Y=c+e$, $D=c+dY$ and use the same product with $X=\max(d+1,Y)$; $e=1$ follows by $(b,c,d,e)\mapsto(e,d,c,b)$.
For $d=1$, expand
\[
 m-2p=2+b^2+b+c^2+e^2+ce-2be+2c-f(D-f+2)
\]
and use $(f-X)(D-f-X)\ge0$ with $X=\max(b+2,c+e+1)$; $c=1$ is the same symmetry. All rows contradict the MA endpoint bounds. Hence the guard holds, proving the bounded routes.
\end{proof}

In a word minimizing high-state count and then length, the initial alternating unit block before a third pivot has at most three high reached states, by this lemma and \eqref{eq:tailtransport}. Thus it is $01$ or $010$. To exclude $010$, the pentagon gives
$\mut{1,0}C=(\mut{0,1,0}C)^{(1,0,2,3)}$. If $\mu_1C$ were under the full ceiling $F$, this replacement would remove a high state. The following calculation rules out the required escape.

All its incident entries are bounded by the third-state ceiling, so an escape at $\mu_1C$ can occur only in entry $23$. Exchange $2,3$ to orient the escape negatively. The first and third high states force
\[
 C=(1,-b,c,d,-E,f),\quad b,c,d,E\ge1,\quad d<b,\quad f=F-bc.
\]
The successive states at $0,01,010$ are
\[
 (-1,b,-c,d-b,-E,F),\quad(1,d,-c-E,b-d,E,F),\quad
 (-1,-d,c+E,b-d,-c,F-d(c+E)).
\]
The exact guards reduce to
\[
 bc\le2F-1,\quad2F<bc+dE,\quad c+E\le F,\quad
 c+E=F\ \hbox{or}\ d(c+E)=2F.
\]
The proper scalar and Pfaffian are
\[
 m=1+b^2+c^2+d^2+E^2+cE-bd+F^2-F(bc+dE),\quad p=F-b(c+E)+cd.
\]
If $c+E=F$, the guards force $b\ge3$ and $K=2b^2-3b-4d\ge1$, and
\[
 b\{2(-p)-F\}=FK+2d(2F-bc)>0.
\]
Otherwise set $L=c+E$, $z=b-d$. Then $dL=2F$, $3\le d,L\le F-1$, $F\ge5$, and
\[
 m=1+d^2+L^2-cE-F^2+z^2+dz-Fcz\le d^2+L^2-cE-F^2.
\]
The nonnegative products $(d-3)(F-1-d)$, $(L-3)(F-1-L)$ and $(d-3)(L-3)$ give $3(d^2+L^2)\le2F^2-5F+36$. Hence $3m\le-F^2-5F+36-3cE<0$. Both alternatives contradict \eqref{eq:prefixscalar}. The first third pivot therefore follows exactly $01$.

\subsection{The exhaustive three-pivot sign split}\label{app:signs}
Write the low lead as $C=(1,x,c,d,e,f)$. The first two reached states are high and the first three states, at $0,01,012$, are below $F$. There are precisely three selected-triangle sign cases.

\subsubsection{Cyclic signs: $x=-b<0<d$}
If $c>0$, the first crossing is $f+bc=F$. If $e\ge0,b<d$, the Pfaffian $F+c(d-b)+be$ exceeds its budget. If $e\ge0,d<b$, the second $23$ entry $F+(b-d)e$ forces $e=0$. For $e=-E<0,d<b$, the third $13$ entry $E+(b-d)F$ forces $b-d=1,E=0$. In these smaller-side cases
\[
 m=1+b^2+c^2+d^2+f^2-bd+bcf,\qquad p=F-c.
\]
Now $2(F-c)\le F-1$ and $bc\le2F-1$ imply $b\le3$. If $b=3,d=2$, put $v=F-c$; then $v\ge2$, $c\ge v+1$, and $m=8+v^2-cF<0$. Thus $b=2,d=1$, giving $A_c$.

The remaining case $c>0,e=-E<0,b<d$ uses the second-high hypothesis. Set $r=d-b\ge1$. The first states are
\[
 (-1,b,-c,r,-E,F),\quad(1,b,-c-E,-r,E,F-rE).
\]
The second wall is $c+E=F$ or $rE=2F$. The third $13$ bound excludes the latter. On $c+E=F$,
\[
 m=1+b^2+br+r^2-(b-2)F^2-rFE-cE,\quad p=(r+1)F-(b+r)E.
\]
The inequalities $3\le m\le F-1$, $2|p|\le m$, $b+r\le F-1$, $rE\le2F$ force $b=r=1$. Indeed $b\ge3$ gives $m<3$; at $b=2$ use
\[
 m=6+2r+r^2-(r+1)F-(E-1)\{(r+1)F-E-1\}<3.
\]
At $b=1$, $p=(r+1)c$ yields $E\ge(2r+1)c+1$, so $r\le2$; for $r=2$, $m=8+c^2-EF<3$. The last matrix, with $r=1$, becomes $I_{c,c}$ after a signed swap of $1,2$, and is excluded below.

If $c\le0$, the first mutation is $(-1,b,-c,d-b,e+c,f)$. A crossing requires $c=-t<0$, $e=-v<0$, $t+v=F$. Existence of a proper class forces $f\ge0$, since the all-negative fourth-row chamber has none. If $d<b$, the third $13$ entry $F+(b-d)f$ forces $f=0$, violating the Pfaffian bound. If $d=b$, that bound forces $b=1$ and
$m=2+t^2+v^2+f^2+tv-vf>t+v=F$.
If $d>b$, the third $23$ entry $(d-b)F-f$ forces $d-b=1$, and
\[
 bv+td-f=(b-1)F+t+(F-f)
\]
forces $b=1,d=2$ and $1\le F-f\le T$. With $u=F-f$ this is $I_{t,u}$.

The remaining equal-connector wall $b=d$, $c>0$ is independent of these strict comparisons. The first two tuples are
\[
 (-1,b,-c,0,e-[c]_-,f+b[c]_+),\quad
 (1,b,-c-[e-[c]_-]_-,0,-e+[c]_-,f+b[c]_+).
\]
If $c=0$ there is no crossing. For $c>0$ the Pfaffian bound forces $e=-E<0$, while the third ceiling gives $bF-c-E\le F$, hence $b\le2$. At $b=1$,
$m=2+c^2+E^2+f^2+cf-Ef>c+f=F$.
At $b=2$, set $v=F-c$, $z=F-E$. Then $v,z\ge1$, $v+z\le F$, and
$m=5+v^2+z^2+vz-F(v+z)$. The condition $m\ge3$ forces $vz\le2$, $v+z\le3$; with $F\ge4$ this gives $m<3$, a contradiction. This completes the cyclic reduction.

\subsubsection{Left signs: $x=-b<0$, $d=-D\le0$}
Set $h=e-[c]_-$, $g=f+b[c]_+$, $z=g+(b+D)[h]_+$. The first states and two decisive third entries are
\begin{align*}
 P&=(-1,b,-c,-D-b,h,g),\\
 Q&=(1,-D,-c-[h]_-,D+b,-h,z),\\
 (\mu_2Q)_{01}&=1-D(D+b),\qquad
 (\mu_2Q)_{13}=-h+(D+b)[z]_+.
\end{align*}
A first crossing lies on $b+D=F$, $h=-F$, or $g=F$.
On $h=-F$, write $c=-t,e=-u$, $t,u\ge1$, $t+u=F$; the third ceiling forces $f\le0$. The proper scalar is
\[
 1+b^2+D^2+t^2+u^2+f^2+tu+bD-f\{(b+D)u+bt\}>F.
\]
On $g=F$, the second and third ceilings force $h=0,b=1,D=0,e=0,c,f\ge1,c+f=F$. The scalar $2+c^2+f^2+cf$ is greater than $F$.
On $b+D=F$, the third $01$ bound gives $D=1,b=T$, and the third $13$ bound gives $z\le1$. If $c>0$, then $e\le0$ and $P=f+Te-c\le-T$, contrary to $2|P|\le T$. If $c\le0<f$, then $f=1,c=e=0$ and $m=T^2+T+3>T$. Otherwise put $t=-c\ge0$; the Pfaffian bound gives $e\le1$ and
\[
 m=T^2+T+2+t^2+e^2-te+f^2+f\{(T+1)e-Tt\}>T.
\]
Use $t^2+e^2-te\ge0$, $(T+1)e-Tt\le T+1$, $f\le0$ and completion of the square. Thus every left-sign wall contradicts the budget.

\subsubsection{Right signs: $x\ge0$}
A crossing requires $c=-t<0$, $1\le t\le T$. Put $z=f-xt+\Delta(-d,e-t)$. Then
\begin{align*}
 P&=(-1,-x,t,d,e-t,f-xt),\\
 Q&=(1,-x-[d]_-,t-[e-t]_-,-d,t-e,z),\\
 (\mu_2Q)_{01}&=1-(x+[d]_-)[d]_-,\\
 (\mu_2Q)_{03}&=t-[e-t]_--(x+[d]_-)[z]_-,\\
 (\mu_2Q)_{13}&=t-e+\Delta(-d,z).
\end{align*}
The first wall is $e-t=-F$ or $f-xt=-F$; the Pfaffian is $f-xe-td$.
On the incident wall set $u=F-t\ge1$, so $e=-u$. If $d<0$, the third $13$ bound gives $f-xt\le0$ and the $03$ bound gives $xt-f\le t$; the Pfaffian bound forces $x=0$. If $d\ge0$, the correction is $-dF$, giving $d\le1$. At $d=0$ the remaining bounds give $x\le1$, with $f\ge0$ when $x=1$; at $d=1$ they give $x=0,f\ge0$. The exhaustive scalar list is
\begin{center}\begin{tabular}{ll}\toprule
$(x,d)$ and additional guard & $m$\\\midrule
$(0,-D)$, $D\ge1$, $f\le0$ & $1+D^2+t^2+u^2+f^2+tu-uDf$\\
$(0,0)$ & $1+t^2+u^2+f^2+tu$\\
$(0,1)$, $f\ge0$ & $2+t^2+u^2+f^2+tu-(t+u)f$\\
$(1,0)$, $f\ge0$ & $2+t^2+u^2+f^2+tu-tf$\\\bottomrule
\end{tabular}\end{center}
Every row exceeds $F=t+u$, by nonnegativity in the first two and completion of squares in the last two, using $t,u\ge1$.

On the opposite wall $f-xt=-F$, the second-high condition is essential. For $d<0$ the possible second walls are $x-d=F$ or $z=\pm F$. The third $01,03,13$ bounds exclude each: the first forces $d=-1$ and then violates $e\le F-1$; $z=-F$ forces $e=t$ and violates the $03$ bound; $z=F$ forces $d=-1,e-t=2F$. If $d\ge0$, the third $03$ bound gives $x=1,e\ge0$. If $d>0$, the second ceiling gives $[e-t]_-=0$ and the third $13$ bound then gives $d=1,e=t$, whose Pfaffian is too large. Hence $d=0$. Put $u=F-t$, so $f=-u$; the scalar
\[
 m=2+t^2+u^2+e^2+tu-te>t+u=F
\]
again contradicts the budget. This finishes all selected-triangle signs, including all zero boundaries.

\subsection{The two surviving shapes and the equal wall}\label{app:equal}
For $A_c=(1,-2,c,1,0,F-2c)$, $1\le c\le T$, the relevant reached tuples are
\begin{center}\small\begin{tabular}{c|rrrrrr}\toprule
Word & $01$&$02$&$03$&$12$&$13$&$23$\\\midrule
$0$&$-1$&$2$&$-c$&$-1$&$0$&$F$\\
$01$&$1$&$1$&$-c$&$1$&$0$&$F$\\
$012$&$1$&$-1$&$F-c$&$-1$&$F$&$-F$\\
$0120$&$-1$&$1$&$c-F$&$-2$&$F$&$-c$\\
$01202$&$1$&$-1$&$c-F$&$2$&$F-2c$&$c$\\\midrule
$1$&$-1$&$-1$&$c$&$-1$&$0$&$F-2c$\\
$10$&$1$&$1$&$-c$&$-1$&$c$&$F-c$\\
$102$&$2$&$-1$&$F-2c$&$1$&$c$&$c-F$\\\bottomrule
\end{tabular}\end{center}
At $012$, the nonreturn pivots $1,3$ create $2F-c>F$ and $F^2-1>F$ respectively, forcing $0$. At $0120$, pivots $1,3$ create $2F-c>F$ and $-1-F(F-c)<-F$, forcing $2$. Thus the first return is exactly $01202$. All three replacement states at $102$ are low, and for $p=(1,2,0,3)$,
\[
 \mut{1,0,2}A_c=(\mut{0,1,2,0,2}A_c)^p.
\]
Four high states are removed from the whole word by \eqref{eq:tailtransport}.

For $I_{t,u}=(1,-1,-t,2,t-F,F-u)$ the states at $0,01,012$ are
\[
 (-1,1,t,1,-F,F-u),\quad(1,1,t-F,-1,F,-u),\quad(2,-1,t-F,1,F-u,u).
\]
Exactly the first two are high. Both replacement words terminate at
$E=(1,-2,F-u,1,u,t-F)=(\mut{0,1,2}I_{t,u})^{(1,2,0,3)}$.
For $t>u$, the first four states of $10212$ are
\begin{align*}
 &(-1,1,-t,-2,F-t,2t-F-u),\quad(1,-1,t,-1,F-t,t-F-u),\\
 &(1,1,2t-F-u,1,-u,F+u-t),\quad(-1,2,2t-F-u,-1,u,F-t).
\end{align*}
For $u>t$, the first four states of $20210$ are
\begin{align*}
 &(-1,1,-t,-2,F+t-2u,u-F),\quad(1,-1,t,-1,F+t-2u,u-F-t),\\
 &(1,1,u-F,1,-u,F+t-u),\quad(-1,2,u-F,-1,u,F+t-2u).
\end{align*}
All entries in the appropriate strict range, including those of $E$, have absolute value at most $T$. For example $F+u-t\le T$ exactly when $t\ge u+1$; the other entries follow from $1\le t,u\le T$. Each replacement removes two high states despite adding two letters.

It remains to exclude $t=u$. Every proper class of $I_{t,u}$ has scalar $4+t^2-tu+u^2$. Consequently $t=u$ gives $m=4+t^2$ and $F\ge t^2+5$.
If $t=1$, then $(m,P^2)=(5,4)$. A topological acyclic endpoint has entries in $\{0,1,2\}$. Checking these $3^6$ possibilities gives, up to signed permutation, a member of the following closed core:
\begin{align*}
 K_0&=(-2,0,1,0,-1,-1),& K_1&=(-1,-1,0,-1,1,-1),\\
 K_2&=(-1,-1,0,0,-1,1),& K_3&=(-2,0,0,0,0,-1).
\end{align*}
The complete mutation table is
\begin{center}\begin{tabular}{c|cccc}\toprule
&$0$&$1$&$2$&$3$\\\midrule
$K_0$&$K_0^{1023}$&$K_0^{1023}$&$-K_0^{1023}$&$-K_1^{2031}$\\
$K_1$&$-K_1^{0213}$&$-K_0^{1302}$&$-K_1^{1023}$&$K_2$\\
$K_2$&$-K_2^{0213}$&$-K_1^{2301}$&$K_2^{1302}$&$K_1$\\
$K_3$&$-K_3^{0132}$&$-K_3^{0132}$&$K_3^{0132}$&$K_3^{0132}$\\\bottomrule
\end{tabular}\end{center}
Here the four digits denote the new-to-old permutation. Undoing an acyclicizing word puts the original matrix in this height-two core. But $I_{1,1}$ has an entry of magnitude $F-1\ge5$, a contradiction. This small endpoint enumeration and the displayed transition table are finite computer-assisted components of the height argument, independently replayed in the supplied checker.

For $t\ge2$, write $q=F-t\ge t^2-t+5\ge7$ and set
\[
 C_0=(1,-1,-t,2,-q,q),\qquad C_1=(-1,1,t,1,-F,q).
\]
Both are cyclic. Their first three mutations are
\begin{center}\begin{tabular}{c|ccc}
&$0$&$1$&$2$\\\hline
$C_0$&$C_1$&$C_0^{0213}$&$C_0^{0213}$\\
$C_1$&$C_0$&$-C_1^{2103}$&$-C_1^{1023}$
\end{tabular}\end{center}
The remaining mutations are
\[
 \mu_3C_0=R_0^{2310},\quad \mu_3C_1=-R_1^{3120},
\]
where
\[
 R_0=(q,-t,-q,1+tq,q^2-2,1),\quad
 R_1=(F,-q,-t,Fq-1,tF-1,1).
\]
Both belong to the following unbounded root domain:
\begin{equation}\label{eq:rootdomain}
 R=(A,-D,-E,B,C,1),\quad A,D,E\ge2,\quad B>\max(A,D),\quad C>\max(A,E).
\end{equation}
Here the return label is $0$ and the special pair is $23$. Put
$h=BD-A$, $j=CE-A$, $\ell=B(D+E)+CE-A$, and define
\begin{align*}
 W&=(-h,D,-E,-B,B+C,-1),& U&=(-\ell,D,E,C,-B-C,1),\\
 U'&=(-\ell,-D,D+E,-C,-B,-1),& W'&=(-j,E,-D-E,-C,B,1).
\end{align*}
Together with their $23$ swaps these are a ten-state cyclic cell, with exact chain
$R\xrightarrow{2}W\xrightarrow{3}U\xrightarrow{2}U'\xrightarrow{3}W'\xrightarrow{2}R^{0132}$.
All remaining transition types are
\begin{center}\begin{tabular}{c|cccc}\toprule
&$0$&$1$&$2$&$3$\\\midrule
$R$&parent&new root&$W$&$(W')^{0132}$\\
$W$&fork&fork&$R$&$U$\\
$U$&new root&fork&$U'$&$W$\\
$U'$&fork&new root&$U$&$W'$\\
$W'$&fork&fork&$R^{0132}$&$U'$\\\bottomrule
\end{tabular}\end{center}
The other five rows are their $23$ swaps. After signed relabelling the new-root parameter lists $(A',B',C',D',E')$ are
\[
 (A,AB-D,AC-E,B,C),\quad
 (\ell,\ell D+C,\ell E-B-C,D,E),\quad
 (\ell,D+\ell C,\ell B-D-E,C,B).
\]
For these three rows respectively, the mutated matrix is the new root relabelled by $(1,0,2,3)$, the new root itself, and the negative of the new root relabelled by $(1,0,2,3)$. They satisfy \eqref{eq:rootdomain}, because $h>\max(B,D)$, $j>\max(C,E)$ and $\ell>\max(B+C,D+E)$ follow from that domain. A normalized fork with return $0$ is
$(-a,b,c,-d,-e,f)$, where $a,b,c,f\ge2$, $d>\max(a,b)$ and $e>\max(a,c)$. Direct substitution shows that each fork entry of the table has this form in a signed frame. For completeness, the signed frames taking the six fork exits to this displayed form are
\[
\begin{array}{c|rrrrrr}
\text{exit}&\mu_0W&\mu_1W&\mu_1U&\mu_0U'&\mu_0W'&\mu_1W'\\
\text{sign}&+&-&-&+&+&-\\
\text{frame}&0213&1203&1320&0321&0213&1203
\end{array}
\]
Thus the claimed inequalities can be checked by substituting $h,j,\ell$ and the root-domain inequalities, without searching for a frame. Mutation at any nonreturn vertex stays a cyclic fork, by \cite[Lemma~2.5]{Warkentin14}; mutation at the return vertex undoes the entry.

Start with the signed-permutation copies of $C_0,C_1$ and close under attaching these root cells and fork branches to included parents. Every state is cyclic. The tables and fork continuation prove closure under every mutation: a transition stays in a cell, returns to its parent, or attaches a new cell or fork. Induction on a finite attachment construction, and then on word length, includes every descendant of $C_0$. Hence $C_0=I_{t,t}$ is NMA, excluding the equal wall without a bound on $t$. Combined with the exhaustive sign split and the whole-word replacements, this completes the height proof.
